% SOURCE_STATEMENT: "Given a Lorentz gas in the Boltzmann–Grad limit with initial distribution $f_0(x,v)$, does $f^\\varepsilon(x,v,t)$ converge as $\\varepsilon\\to0$ to a limit $f(x,v,t)$, and, if so, what equation does $f$ satisfy?"
% Preparation and final self-audit: CONTRIBUTING.md prompts C and D; source/audit artifacts retained in the local physics_statements_2026_10_11 batch.
% !TeX program = lualatex
\documentclass[11pt]{article}
\usepackage[margin=25mm]{geometry}
\usepackage{fontspec}
\setmainfont{Latin Modern Roman}
\usepackage{amsmath,amssymb,mathtools}
\usepackage{unicode-math}
\setmathfont{Latin Modern Math}
\usepackage{xurl,hyperref}
\hypersetup{colorlinks=true,urlcolor=blue}
\setcounter{secnumdepth}{0}
\setlength{\parindent}{0pt}
\setlength{\parskip}{5pt}
\setlength{\emergencystretch}{3em}
\newcommand{\sref}[2]{\hyperlink{source:#1:#2}{[S#2]}}
\newcommand{\cataloguescope}[1]{\paragraph{Recorded target.}#1}
\newcommand{\R}{\mathbb R}
\newcommand{\C}{\mathbb C}
\newcommand{\N}{\mathbb N}
\newcommand{\Z}{\mathbb Z}
\newcommand{\Q}{\mathbb Q}
\newcommand{\F}{\mathbb F}
\newcommand{\eps}{\varepsilon}
\begin{document}
% UnsolvedMath ID 30001644; source catalogue code OWR-4536-006
\setcounter{section}{30001643}
\section{Boltzmann–Grad Limits for the Lorentz Gas}\label{30001644}
{\small UnsolvedMath ID 30001644 · Mathematical Physics}
\subsection{Definitions and mathematical statement}

\paragraph{Shared notation.}$\mathbb N=\{1,2,\ldots\}$; the standard number fields have their usual meanings. All additional conventions are stated locally.


\paragraph{Definitions and conventions.} A Lorentz gas consists of point particles moving freely at fixed speed between elastic scatterers and reflecting specularly at collisions. The particle density is $f^\varepsilon(x,v,t)$ with $x\in\mathbb R^n$ and $v\in S^{n-1}$. The Boltzmann--Grad scaling shrinks scatterers and rescales their number or spacing to retain the collision scale. Besides random and periodic arrays, the report examines aperiodic cut-and-project quasicrystal centers. Convergence requires a topology on the particle distributions and an initial density excluding obstacle interiors.

\paragraph{Source scope and qualifications.} The context is quasicrystal scatterers, which the short catalogue question omits; random and periodic arrays are comparison cases. The report uses a particular radius/spacing normalization and presents simulations rather than proof for its aperiodic model. The limit must not silently be assumed to satisfy the ordinary linear Boltzmann equation.

\paragraph{Catalogue mathematical target.}
\paragraph{Catalogue wording retained for traceability.}
Given a Lorentz gas in the Boltzmann–Grad limit with initial distribution $f_0(x,v)$, does $f^\varepsilon(x,v,t)$ converge as $\varepsilon\to0$ to a limit $f(x,v,t)$, and, if so, what equation does $f$ satisfy?

\paragraph{Saved literature context.} The supplied catalogue's assessment dated 2026-08-21 records: Random and periodic Lorentz gases have rigorous limiting theories, but the requested convergence and equation for quasicrystal scatterers remains unresolved; the report only gives simulations. The assessment is retained as dated provenance; it is not a new comprehensive literature or open-status certification.
\subsection{Short English statement}
Derive the limiting particle dynamics of the quasicrystal Lorentz gas and determine its kinetic equation.
\subsection{Sources}
\begin{itemize}
\item[\hypertarget{source:30001644:1}{S1}] ULAM.AI and the UnsolvedMath Contributors, UnsolvedMath, record 30001644 (OWR-4536-006). Catalogue curation under CC BY 4.0; definitions and English exposition added locally.
\url{https://huggingface.co/datasets/ulamai/UnsolvedMath}.
\item[\hypertarget{source:30001644:2}{S2}] Classical and Quantum Mechanical Models of Many-Particle Systems (2010), Oberwolfach Reports, DOI 10.4171/owr/2010/54, pp. 3159--3236; official report source and mathematical context.
\url{https://doi.org/10.4171/owr/2010/54}.
\item[\hypertarget{source:30001644:3}{S3}] Official EMS Press PDF of the same Oberwolfach report; PDF pages 71--74, printed pages 3229--3232.
\url{https://ems.press/content/serial-article-files/46313}.
\item[\hypertarget{source:30001644:4}{S4}] kinetic theory of the low-density Lorentz gas; reference in the saved dated catalogue literature review.
\url{https://people.maths.bris.ac.uk/~majm/bib/kinetic-lorentz-MAMS.pdf}.
\end{itemize}
\label{end:30001644}
\end{document}
